\documentclass[10pt,a4paper]{amsart}
\usepackage[margin=19mm]{geometry}
\usepackage{amsmath,amssymb,mathtools}
\usepackage[scr=rsfs,cal=euler]{mathalfa}
\usepackage{xfrac}
\usepackage[unicode,hidelinks,bookmarks=false]{hyperref}
\newcommand{\ie}{i.e.\ }
\newcommand{\cf}{cf.\ }
\newcommand{\resp}{respectively\ }
\newcommand{\Subsection}{\subsection}
\providecommand{\nobreakdash}{\nobreak\hbox{-}}
\setlength{\emergencystretch}{2em}
\multlinegap0pt
\usepackage{amssymb}
\usepackage{mathtools}
\newtheorem{lem}{Lemma}
\title{On the Minimax Bifurcation Formula}
\author{Y. Sh. Il'yasov}
\date{Source: arXiv:2605.17331v2 (CC BY 4.0)}
\begin{document}
\maketitle

\noindent\emph{Source and licence.} On the Minimax Bifurcation Formula. Version arXiv:2605.17331v2 (CC BY 4.0), distributed by arXiv under the Creative Commons Attribution 4.0 International licence, http://creativecommons.org/licenses/by/4.0/, as recorded in the arXiv licence metadata for this exact version and shown on the abstract page https://arxiv.org/abs/2605.17331v2. Full licence notice: \texttt{licenses/arxiv-2605.17331-CC-BY-4.0.txt}.\par\smallskip

\noindent\textit{Editorial scope.} Counted unit: the article's lem lem:PS-application (source0.tex lines 1757--1762). The article's complete original proof of it is delivered with it (source0.tex lines 1764--1992, one \texttt{proof} environment of 229 lines). No mathematics is written, repaired or re-derived: the statement and the proof are the article's own lines, in the article's own notation, and no retained line is substituted or re-flowed.  Retained uncounted as the source-local context the proof needs: lines 1507--1512; lines 2871--2881.  Nothing was omitted from any retained span, so the package carries no omission record.  No retained line contains a citation, so the wrapper carries no bibliography and no bibliography entry.  No retained line contains a figure environment, an image-inclusion command or drawing code, so no image is delivered and the package declares no asset dependency.  Editorial formatting only, in addition: an AMS \texttt{amsart} preamble that restates the article's own declarations and macros used by the retained lines, namely the environments \texttt{lem} on separate counters, together with the standard packages those lines need.  The retained environments print in this document as Lemma 1 in order of appearance, whereas the article prints the same items with the numbering of its own full document.  The complete native source of the article as distributed in the arXiv e-print for this exact version is bundled unchanged as \texttt{sources/200786/source0.tex}.
\begin{equation}
	\label{eq:picone-superlinearity-f}
	Q_f(x,u)\ge \theta\sum_{i=1}^m f_i(x,u)u^i
	\qquad
	\text{for a.e. }x\in\Omega,\quad u\in\mathbb R_+^m,
\end{equation}

\begin{lem}[Discrete comparison for a sublinear problem]
	\label{lem:discrete-comparison-sublinear}
	Let \(A_r\) be a nonsingular \(M\)-matrix, \(0<p<1\), and \(c_0>0\).
	Suppose that \(U,W\in S_r^o\), written in nodal coordinates, satisfy
	\[
	A_rU\ge c_0U^p,
	\qquad
	A_rW=c_0W^p ,
	\]
	where the powers are understood componentwise. Then \(U\ge W\).
\end{lem}

\begin{lem}\label{lem:PS-application}
	Under the assumptions stated above, suppose in addition that the discrete
	comparison estimate of Lemma~\ref{lem:discrete-comparison-sublinear}
	holds for the Galerkin operator \(A_r\). Then \(\mathcal R\) satisfies the
	extended \((PS)_e\)-condition at every level \(\lambda>0\).
\end{lem}

\begin{proof}
	Let \(r_n\to\infty\), and let
	\((u_n,v_n)\in\mathbb S_{r_n}^o\times\mathbb S_{r_n}^o\) be an extended
	\((PS)_e\)-sequence at a level \(\lambda>0\). Set
	\[
	\lambda_n:=\mathcal R(u_n,v_n)\to\lambda .
	\]
	By the homogeneity of \(\mathcal R\) with respect to the second variable,
	we normalize \(a_m(v_n,v_n)=1\). The Galerkin criticality relations,
	equivalent to \(D_u\mathcal R(u_n,v_n)=D_v\mathcal R(u_n,v_n)=0\) on
	\(\mathbb W_{r_n}\), are
	\[
	\langle\mathcal F(u_n,\lambda_n),\zeta\rangle=0
	\quad \forall \zeta\in\mathbb W_{r_n},
	\qquad
	\bigl\langle D_u\mathcal F(u_n,\lambda_n)\xi,v_n\bigr\rangle=0
	\quad \forall \xi\in\mathbb W_{r_n}.
	\]
	Testing the first identity by \(u_n\), we get
	\begin{equation}\label{eq:energy-ps-g}
		a_m(u_n,u_n)
		=
		\lambda_n\langle g(u_n),u_n\rangle_m
		+
		\langle f(u_n),u_n\rangle_m .
	\end{equation}
	
	We first prove that \((u_n)\) is bounded in \(\mathbb W\). Let
	\(\theta_n\in\mathbb W_{r_n}\) be the nodal function defined componentwise by
	\[
	(\theta_n)_i^k=\frac{(u_{n,i}^k)^2}{v_{n,i}^k}.
	\]
	Since \(u_n,v_n\in\mathbb S_{r_n}^o\), this function is well defined. Taking
	\(\xi=\theta_n\) in the adjoint Galerkin identity, using the componentwise
	discrete Picone inequality, and applying the arithmetic--geometric mean
	inequality to the cooperative off-diagonal terms, we obtain
	\[
	a_m(u_n,u_n)
	\ge
	(q-1)\lambda_n\langle g(u_n),u_n\rangle_m
	+
	\int_\Omega Q_f(x,u_n)\,dx .
	\]
	By \eqref{eq:picone-superlinearity-f}, for some \(\eta>1\),
	\[
	a_m(u_n,u_n)
	\ge
	(q-1)\lambda_n\langle g(u_n),u_n\rangle_m
	+
	\eta\langle f(u_n),u_n\rangle_m .
	\]
	Combining this estimate with \eqref{eq:energy-ps-g}, we find
	\[
	(\eta-1)a_m(u_n,u_n)
	\le
	(\eta-q+1)\lambda_n\langle g(u_n),u_n\rangle_m .
	\]
	Since
	\[
	\langle g(u_n),u_n\rangle_m
	=
	\|u_n\|_{\mathbb L^q}^q
	\le
	Ca_m(u_n,u_n)^{q/2},
	\]
	and \(q<2\), the sequence \((u_n)\) is bounded in \(\mathbb W\).
	
	By the normalization of \(v_n\) and the coercivity of \(a_m\), the sequence
	\((v_n)\) is also bounded in \(\mathbb W\). Hence, after passing to a
	subsequence,
	\[
	u_n\rightharpoonup u^*,
	\qquad
	v_n\rightharpoonup v^*
	\quad\text{weakly in }\mathbb W,
	\]
	and
	\[
	u_n\to u^*,
	\qquad
	v_n\to v^*
	\quad\text{strongly in }\mathbb L^\beta,\quad 1\le\beta<2^*,
	\]
	and a.e. in \(\Omega\). Since \(\mathbb S\) is weakly closed and
	\(\mathbb S_{r_n}\subset\mathbb S\), we have \(u^*,v^*\in\mathbb S\).
	
	We now pass to the limit in the Galerkin equation. Let
	\(\zeta\in\mathbb W\). Choose \(\zeta_n\in\mathbb W_{r_n}\) such that
	\[
	\zeta_n\to\zeta
	\qquad\text{strongly in }\mathbb W .
	\]
	Then
	\[
	\langle\mathcal F(u_n,\lambda_n),\zeta_n\rangle=0.
	\]
	Using the weak convergence of \(u_n\), the strong convergence in
	subcritical Lebesgue spaces, and the growth assumptions on \(f\), we obtain
	\[
	\langle\mathcal F(u^*,\lambda),\zeta\rangle=0
	\qquad \forall \zeta\in\mathbb W .
	\]
	
	We next show that \(u^*\in\mathbb S^o\). Since
	\(\lambda_n\to\lambda>0\), for all large \(n\),
	\[
	\mathcal L u_n^k
	=
	\lambda_n(u_n^k)^{q-1}+f_k(x,u_n)
	\ge
	\frac{\lambda}{2}(u_n^k)^{q-1}
	\]
	in the discrete weak sense, for \(k=1,\ldots,m\). By
	Lemma~\ref{lem:discrete-comparison-sublinear}, applied componentwise,
	\[
	u_n^k\ge w_n^k,
	\]
	where \(w_n^k\in S_{r_n}^o\) is the Galerkin solution of the scalar
	sublinear problem
	\[
	\mathcal L w_n^k=\frac{\lambda}{2}(w_n^k)^{q-1}.
	\]
	The standard convergence of Galerkin solutions for this scalar sublinear
	problem gives
	\[
	w_n^k\to w
	\quad\text{strongly in }H_0^1(\Omega)
	\quad\text{and a.e. in }\Omega,
	\]
	where \(w\) is the positive weak solution of
	\[
	\mathcal L w=\frac{\lambda}{2}w^{q-1},
	\qquad
	w\in H_0^1(\Omega),
	\qquad
	w>0 \quad\text{in }\Omega .
	\]
	Passing to the limit in \(u_n^k\ge w_n^k\), we obtain
	\[
	u^{*,k}\ge w>0
	\quad\text{a.e. in }\Omega,
	\qquad k=1,\ldots,m .
	\]
	Since \(u^*\) solves the system and the right-hand side has subcritical
	growth, standard bootstrap/Moser estimates and boundary H\"older regularity
	imply \(u^*\in(C(\overline\Omega))^m\). Hence
	\(u^{*,k}>0\) in \(\Omega\), \(k=1,\ldots,m\), and therefore
	\(u^*\in\mathbb S^o\).
	
	It remains to strengthen the convergence of \(u_n\) and to prove that
	\(v^*\ne0\). Testing the limit equation by \(u^*\), we get
	\[
	a_m(u^*,u^*)
	=
	\lambda\langle g(u^*),u^*\rangle_m
	+
	\langle f(u^*),u^*\rangle_m .
	\]
	Comparing this identity with \eqref{eq:energy-ps-g}, and using the compact
	convergence of the nonlinear terms, we obtain
	\[
	a_m(u_n,u_n)\to a_m(u^*,u^*) .
	\]
	Together with the weak convergence \(u_n\rightharpoonup u^*\), the
	coercivity of \(a_m\) gives
	\[
	u_n\to u^*
	\qquad\text{strongly in }\mathbb W .
	\]
	
	Finally, taking \(\xi=v_n\) in the adjoint Galerkin identity and using the
	normalization \(a_m(v_n,v_n)=1\), we obtain
	\[
	1
	=
	(q-1)\lambda_n
	\left\langle u_n^{q-2}v_n,v_n\right\rangle_m
	+
	\left\langle f_u(u_n)v_n,v_n\right\rangle_m .
	\]
	On the other hand, testing the Galerkin equation by the nodal function with
	values
	\[
	\frac{(v_{n,i}^k)^2}{u_{n,i}^k}
	\]
	and using the componentwise discrete Picone inequality yields
	\[
	1
	\ge
	\lambda_n
	\left\langle u_n^{q-2}v_n,v_n\right\rangle_m
	+
	\left\langle f(u_n),\frac{v_n^2}{u_n}\right\rangle_m .
	\]
	Choose \(\chi\in(q-1,1)\). Multiplying the last inequality by \(\chi\) and
	subtracting it from the preceding identity gives
	\[
	1-\chi
	\le
	\left\langle f_u(u_n)v_n,v_n\right\rangle_m ,
	\]
	because the remaining terms are nonpositive:
	\[
	(q-1-\chi)\lambda_n
	\left\langle u_n^{q-2}v_n,v_n\right\rangle_m\le0,
	\qquad
	-\chi
	\left\langle f(u_n),\frac{v_n^2}{u_n}\right\rangle_m\le0 .
	\]
	If \(v^*=0\), then by compactness
	\[
	v_n\to0
	\quad\text{strongly in }\mathbb L^\beta,\qquad 1\le\beta<2^* .
	\]
	Since \((u_n)\) is bounded in \(\mathbb W\) and \(f_u\) has subcritical
	growth, \(\left\langle f_u(u_n)v_n,v_n\right\rangle_m\to0\),
	contradicting the previous estimate. Hence \(v^*\ne0\).
	
	Consequently, after passing to a subsequence,
	\[
	u_n\to u^*
	\quad\text{strongly in }\mathbb W,
	\qquad
	v_n\rightharpoonup v^*
	\quad\text{weakly in }\mathbb W,
	\]
	with \(u^*\in\mathbb S^o\), \(v^*\in\mathbb S\setminus\{0\}\). This is
	precisely the extended \((PS)_e\)-condition at the level \(\lambda>0\).
\end{proof}

\end{document}
